\documentclass[a4paper]{amsart}
\usepackage[a4paper,margin=1in]{geometry}
\usepackage{amsmath,amssymb,amsthm}
\usepackage{booktabs,array}
\usepackage[parfill]{parskip} 

%--- theorem environments ---------------------------------------------------
\newtheorem{theorem}{Theorem}
\newtheorem{corollary}[theorem]{Corollary}
\newtheorem{remark}[theorem]{Remark}
\theoremstyle{definition}
\newtheorem{assumption}{Assumption}

\begin{document}

\title{Mechanistic Fixed-Bed CO2 Adsorption Model: Analysis and Formulation}
\date{}
\maketitle

\makeatletter
\def\subsubsection{\@startsection{subsubsection}{3}%
  \z@{.5\linespacing\@plus.7\linespacing}{.1\linespacing}%
  {\normalfont\itshape}}
\makeatother

%--- Abstract ---------------------------------------------------------------
\begin{abstract}
A single component CO2 adsorption from an inert carrier (N2) in a 1-D fixed-bed is modelled and derived from conservation laws, closed with LDF kinetics and a nonlinear isotherm (Toth, Dual-Site Langmuir). Fitted analytical models (Yoon-Nelson, Thomas, Bohart-Adams, Clark, fractal sigmoids, etc) parameterise the shape of one curve at one operating point, with $k_{YN}=k\,b\,c_f$ and $\tau=t_{st}(u,c_f,L)$. That composition is exactly why fitted $k_{YN},\tau$ cannot extrapolate across $u$, $c_f$, $L$, or $T$, and why the PDE model below can.
\end{abstract}

\section{Assumptions}
\begin{itemize}

\item \textbf{A1 (1-D):} Radial gradients neglected. Caveat for the bench rig:
\[
    \frac{d_{\mathrm{col}}}{d_p} \lesssim 10
\]
invites wall channeling; treat radial uniformity as an idealisation to be checked, not a fact.

\medskip

\item \textbf{A2 (dilute feed / constant \(u\)):} CO$_2$ mole fraction
\[
    y_f \ll 1
\]
so total molar flux, hence \(u\), is \(z\)-independent. At the measured \(10\text{--}15\%\) feeds this is only first-order accurate; Remark~A.1 gives the exact variable-velocity extension and the \(O(y_f)\) error bound.

\medskip

\item \textbf{A3 (ideal gas, isobaric):} \(P\) uniform (pressure drop \(\ll P\));
\[
    c_{\mathrm{tot}} = \frac{P}{RT}.
\]

\medskip

\item \textbf{A4 (Fickian axial dispersion):} All axial mixing mechanisms (molecular + Taylor--Aris + packing) are lumped into a single \(D_L\).

\medskip

\item \textbf{A5 (LDF):} Intraparticle and film resistance are lumped into a single first-order driving force with constant \(k\) (§A.3); valid when one resistance dominates or concentration profiles are near-parabolic.

\medskip

\item \textbf{A6 (pseudo-homogeneous heat):} Gas and solid share one temperature \(T\); valid when interphase exchange is fast,
\[
    \frac{h_f a_p L}{u \rho_g c_{p,g}} \gg 1,
\]
in which case the two-temperature model of \texttt{derivation.md} (§1.3--1.4) collapses onto this one.

\medskip

\item \textbf{A7 (thermodynamic consistency):} The \((-\Delta H)\) in the energy balance equals the isosteric heat implied by the isotherm's temperature dependence (§A.4.3).

\medskip

\item \textbf{A8:} Constant $
    \varepsilon,\;
    \rho_p,\;
    c_{p,s},\;
    \lambda_{\mathrm{eff}},\;
    D_L
$
over the operating window.

\end{itemize}

\section{Governing Equations}

\subsection{Gas-phase CO2 mass balance}\mbox{}\\


\end{document}